%%%%%
%%%%%  Use LUALATEX, not LATEX.
%%%%%
%%%%
\documentclass[]{VUMIFTemplateClass}

\usepackage{indentfirst}
\usepackage{amsmath, amsthm, amssymb, amsfonts}
\usepackage{mathtools}
\usepackage{physics}
\usepackage{graphicx}
\usepackage{verbatim}
\usepackage[hidelinks]{hyperref}
\usepackage{color,algorithm,algorithmic}
\usepackage[nottoc]{tocbibind}
\usepackage{tocloft}

\usepackage{titlesec}
\newcommand{\sectionbreak}{\clearpage}

\usepackage{titlesec}

\setcounter{secnumdepth}{4}
\setcounter{tocdepth}{3}

\titleformat{\paragraph}
{\normalfont\normalsize\bfseries}{\theparagraph}{1em}{}
\titlespacing*{\paragraph}
{0pt}{3.25ex plus 1ex minus .2ex}{1.5ex plus .2ex}

% Create the custom command
\newcommand{\subsubsubsection}[1]{\paragraph{#1}}

\makeatletter
\renewcommand{\fnum@algorithm}{\thealgorithm}
\makeatother
\renewcommand\thealgorithm{\arabic{algorithm} algorithm}

\usepackage{biblatex}
\bibliography{bibliografija}
%% to change the numbering (numeric or alphabetic) of bibliographic sources, make the change in VUMIFTemplateClass.cls, line 139

% Author's MACROS
\newcommand{\EE}{\mathbb{E}\,} % Mean
\newcommand{\ee}{{\mathrm e}}  % nice exponent
\newcommand{\RR}{\mathbb{R}}




\studyprogramme{Software engineering} %Write your study 


% pilnas pavadinimas, aiskina ka visa tai daro
% prasmingas pavadinimas cia
\worktitle{\textit{UniServe}: A Universal Service Management Platform for Reservation and Order-Based Businesses} 
\workauthor{Tadas Riksas, Darius Spruogis, Gustas Mickus, Kajus Bicka}

\supervisor{Kristina Lapin}

\begin{document}
\selectlanguage{english}

\onehalfspacing
\begin{titlepage}
\vskip 20pt
\begin{center}
\includegraphics[scale=0.55]{images/MIF.png}
\end{center}

\makeatletter

\vskip 20pt
\centerline{\bf \large \textbf{VILNIUS UNIVERSITY}}
\vskip 10pt
\centerline{\large \textbf{FACULTY OF MATHEMATICS AND INFORMATICS}}
\vskip 10pt
\centerline{\large \textbf{\MakeUppercase{\@studyprogramme \space study programme}}}

\vskip 80pt
\centerline{\Large \@worktype}
\vskip 20pt
\begin{center}
    {\bf \LARGE \@worktitle}
\end{center}
\begin{center}
    {\bf \Large \@secondworktitle}
\end{center}
\vskip 80pt

\centering{\Large \@workauthor}
\@ifundefined{@secondauthor}{}
{
\vskip 10pt
\centering{\Large \@secondauthor}
}
\vskip 20pt

\centering{
    \begin{tabular}{rcp{.7\textwidth}}
        {\Large Supervisor} & {\Large :} & {\Large \@supervisor}\\[10pt]
        \@ifundefined{@scientificadvisor}{}
            {
                {\Large Scientific advisor} & {\Large :} & {\Large \@scientificadvisor}\\[10pt]
            }
        \@ifundefined{@reviewer}{}
            {
                {\Large Reviewer} & {\Large :} & {\Large \@reviewer}\\[10pt]
            }
    \end{tabular}}


\vskip 110pt

\centerline{\large \textbf{Vilnius}}
\centerline{\large \textbf{\the\year{}}}

\makeatother

\newpage
\end{titlepage}
%\newgeometry{top=2cm,bottom=2cm,right=2cm,left=3cm}
\setcounter{page}{2}


\singlespacing
\selectlanguage{english}
% list of figures, delete if not needed
% \listoffigures 

%list of tables, delete if not needed
% \listoftables

%Turinys
\tableofcontents
\onehalfspacing


% %Section for abbreviations
% \sectionnonum{List of abbreviations} %Leave if necessary
% This section is for when abbreviations are used. For example:

% \begin{tabular}{rcp{.7\textwidth}}
%     {u.d.i.r.v.} & {} & {uniformly distributed independent random variables}
% \end{tabular}

\section{Introduction}
\subsection{Project Title}

\textbf{Full Title:} \textit{UniServe: A Universal Service Management Platform for Reservation and Order-Based Businesses}

\textbf{Short Title:} \textit{UniServe}


\textbf{Explanation:} "Uni" represents \emph{universal}, indicating the platform's flexibility across a variety of business models 
(reservation based businesses such as massage services or order based businesses such as restaurants). 
"Serve" refers to \emph{service}, emphasizing the platform’s role in providing core functionalities for these types of businesses.

\subsection{Problem statement}


% Who is affected by the problem?

% 

% What is the problem?

% 

% Where does this problem occur?

% 

% When does the problem occur?

% 

% Why does the problem occur?

% 

% Why is the problem important?

% 



Traditional methods like handing out menus and taking orders on-site or booking appointments by phone have become slow and inefficient in the last decade. This is especially noticeable during peak business periods such as holidays, weekends, or promotional events. Small to medium-sized businesses in the restaurant and beauty industries — particularly those that do not use modern digital platforms — are most affected. Limited budgets and lack of technical expertise often prevent them from adopting specialized web applications. This creates multiple problems: employees face increased workload and inefficient information exchange, while customers increasingly expect online services like remote ordering and appointment booking. Without accessible digital platforms, 
businesses risk missing potential customers who rely on the internet, leading to reduced customer reach, lower operational efficiency, and decreased customer satisfaction. The lack of automation results in higher recurring costs, prevents the productivity gains that automation can deliver — gains that could allow staff to work more efficiently and focus on higher-value tasks.

% darbo nasumas dideja

\section{Stakeholder analysis}

\subsection{Identifying stakeholders}

\begin{table}[h]
  \centering
  \caption{Stakeholders Identification}
  \begin{tabular}{|c|c|}
    \hline
    Stakeholder Type    & Stakeholders \\ \hline
    Primary             & Employees, Customers \\ \hline
    Secondary           & Client Business Owners \\ \hline
    Indirect (Tertiary) & App Business Owners \\ \hline
    External/Competitors & Square, Toast POS, Restaurant365 \\ \hline
    Technical Support   & Sysadmins, Tech Support \\ \hline
  \end{tabular}
  \label{tab:stakeholders}
\end{table}

% \textbf{Primary Stakeholders} directly and frequently use the system. Client
% business owners (restaurant/salon owners) operate the system daily. Employees
% process orders and manage appointments through the application.

% \textbf{Secondary Stakeholders} are occasionally involved or interact via
% others. Customers use the system when placing orders or booking appointments.
% Managers oversee operations and may need access to system statistics.

% \textbf{Indirect (Tertiary) Stakeholders} influence purchase decisions and
% benefit from the system. App business owners (our company's owners) profit from
% the system's success. Business managers make decisions about adopting the
% system.

% \textbf{Competitors} represent existing market solutions. Users currently
% employing Square, Toast POS, or similar platforms are potential adopters of our
% system.

% \textbf{Development and Maintenance Team} creates and maintains the system.
% Sysadmins handle infrastructure and servers. Tech support staff resolve user
% issues and provide assistance.

% Nielsen‘s principles
% 1. Learnability
% 2. Efficiency of use
% 3. Memorability
% 4. Few and non-
% catastrophic errors
% 5. Satisfaction

\textbf{Primary Stakeholders:}
\begin{itemize}

    \item \textbf{Employees} want an efficient, error-free system that allows them to process orders and book appointments quickly, with an interface that is easy to use throughout the day.
    % memorization - not really important kadangi darbuotojai sistema naudosis visą dieną taigi bet kokiu atveju atsimins
    % primary as customers use system directly
    \item \textbf{Customers} expect an effective, reliable booking and ordering
    experience with quick response times, and accurate information and
    satisfying interactions.
\end{itemize}

\textbf{Secondary Stakeholders:}
\begin{itemize}
    \item \textbf{Client Business Owners} want tools to manage their business,
    track sales, watch over staff, manage stock, and get reports to run their
    business better and reduce costs. They expect and error-free system that is
    reliable and efficient. 
    % this is secondary stakeholder (statistics), indirect (ownder)
\end{itemize}

\textbf{Indirect (Tertiary) Stakeholders:}
\begin{itemize}
    \item \textbf{App Business Owners} want steady income growth, to increase
    market share, and to compete well with other business software companies.
\end{itemize}

\textbf{Competitors:}
\begin{itemize}
% users, specify
    \item \textbf{Square, Toast POS, Restaurant365}. Users currently relying on these platforms (or contemplating using them) expect similar or improved features, affordable pricing, and intuitive interfaces that allow for quick onboarding without extensive technical training.
\end{itemize}

\textbf{Development and Maintenance Team:}
\begin{itemize}
    \item \textbf{Sysadmins} expect reliable system architecture, scalable
    infrastructure, and clear deployment procedures.
    \item \textbf{Tech Support} expect comprehensive documentation, effective
    troubleshooting tools, and manageable user issue resolution processes.
\end{itemize}


% <The name of the first user group> research {for each user group in both primary and secondary stakeholders}
\subsection{Employee research}

\subsubsection{Current users' activities}

\textbf{Café Barista Managing Orders}

Marius, 21, works as a barista in a busy café. He takes orders at the counter and writes them on sticky notes for the kitchen. During the morning rush, notes sometimes fall to the floor or get smudged by spilled drinks. This leads to wrong orders being prepared, which frustrates customers and wastes ingredients. Marius tries to keep the counter tidy, but the lack of a structured system makes mistakes inevitable and thus Marius spends extra time apologizing and remaking orders, which slows down the queue.

\textbf{Salon Receptionist Controlling Bookings}

Lucille, in her twenties, is the receptionist at a busy, well‑known beauty salon. Throughout the day she fields a constant stream of phone calls to arrange appointments. Each booking is first written by hand in a paper diary, then later re‑entered into a spreadsheet. On busy days, at least one record is often missed or entered incorrectly, resulting in double bookings or missed appointments. When a client calls to reschedule, Lucille leafs through pages of the diary to locate the original booking. This takes time, breaks her concentration, and can leave callers waiting on the line. The salon offers no online booking facility, so all scheduling relies entirely on Lucille being available to answer the phone.

\textbf{Restaurant Manager Handling Discounts}

Youmna, in her early twenties, manages a vibrant Lebanese restaurant on a busy high street. To attract customers during quieter weekday afternoons, the restaurant runs a “Taste of Lebanon” promotion, offering discounts on selected platters and traditional dishes.
When diners ask about the offer, Youmna or the waiting staff consult a printed list of eligible dishes and manually work out the reduced price at the till. At busy times, this slows service. Occasionally, the discount is forgotten, prompting awkward conversations and amended bills; at other times, it is applied in error, reducing profit margins.
To monitor the promotion, Youmna records discounted orders in a small ledger. However, during the rush, staff often forget to make entries, leaving the records incomplete. This makes it difficult to assess the effectiveness of the offer or refine it for future campaigns.

\subsection{Customer research}

\subsubsection{Current users' activities}

% 4 level section should be working now
% \subsubsubsection{test test}

\textbf{Looking for a Quick Lunch}

Annie, a 20‑year‑old university student, has a short break between lectures and only 30 minutes to eat. She visits a small café she has not tried before, hoping to find something quick. The menu gives no indication of preparation times, so she chooses a dish at random and waits nearly 10 minutes for a waiter to take her order. She is then told it will take 25 minutes to prepare. With little time left, Annie leaves and buys a takeaway wrap from a nearby chain instead.

\newpage

\textbf{Booking a Haircut}

Felix, in his late forties, needs a haircut before a family gathering. He calls his usual barber, but the line is engaged. After several attempts, he finally gets through, only to learn that his preferred time slot has already been taken. The receptionist offers alternative appointments, but Felix is at work and cannot check his calendar straight away. Frustrated, he decides to book with another salon that can fit him in immediately.

\textbf{Redeeming a Coffee Promotion}

Sarah‑Louise, 27, spots a “Buy One, Get One Free” coffee offer on a café’s social media page. She visits the café to take advantage of the deal, but the barista is unaware of the promotion and has to confirm it with the manager. This causes a delay at the counter and creates an awkward queue behind her. The discount is eventually applied, but Sarah‑Louise leaves feeling the process was poorly organised.

\subsection{Client Business Owners Research (Secondary User Group)}

\subsubsection{Current users' activities}

\textbf{Restaurant Owner Handling Customer Feedback}

Apolline, 39, owns a family restaurant in a popular tourist area. Customers occasionally leave reviews on Google Maps or TripAdvisor, but she only checks them sporadically. Complaints (such as “long wait for a table” or “incorrect order”) can go unnoticed for several days. Apolline also keeps a suggestion box in the restaurant, though very few customers make use of it. She would prefer to gather feedback in a more systematic way and respond swiftly in order to improve the service. Without structured data, it is difficult for her to identify recurring issues (for example, slow service during specific time frames) and to train staff appropriately.

\textbf{Salon Owner Tracking Staff Performance}

Nils, in his forties, runs a busy salon with six stylists. He tries to track performance and service quality, but still relies on a paper diary and casual feedback. Calculating wages and commissions means hours spent combing through notes and receipts. When figures are challenged, he has no firm record to refer to, which not only heightens administrative burden but also sows discontent amongst the team.


% \subsection{Primary Stakeholders Needs}


% Now that we know what are our primary and secondary stakeholders expectations,
% we can take a look into what are their needs in more detail. Specifically into
% how their needs will influence the design of our system.


% \subsubsection{Client Business Owners Needs}

% The main needs of the client business owners are to have a system where they 
% can manage their business operations daily, track sales, manage inventory,
% generate reports (taxes, profits, losses, etc.), and hire or fire employees.
% This need is driven by the fact that business owners need to have a clear overview of their
% business operations to make informed decisions and ensure the success of their
% business.



% The main needs of the client business owners are to have a system that is
% \textbf{error-free}, \textbf{reliable} and \textbf{efficient}. Business owners will need to use the system daily to manage their employees,
% track sales, manage inventory, and generate reports (taxes, profits, losses,
% etc.). Therefore the system must be error-free to ensure that critical data like
% report generation is not compromised. The system must also be reliable to ensure
% that business operations are not disrupted and that business owners can trust
% the system to perform as expected. Finally, the system must be efficient to
% ensure that business owners can quickly and easily perform their tasks without
% wasting time or effort for looking for features or navigating through complex
% menus.

% \subsubsection{Employees Needs}

% The main need of the employees is to have a system that helps them process orders,
% book appointments, and manage customer interactions. This need is driven by the
% fact that employees need to be able to perform their tasks quickly and accurately
% to ensure customer satisfaction and the smooth operation of the business.

% The main needs of the employees are to have a system that is \textbf{easy to
% learn}, \textbf{efficient} and \textbf{error-free}. Employees will need to use
% the system daily to process orders, book appointments, and manage customer
% interactions. Therefore the system must be easy to learn to ensure that
% employees can quickly adapt to the system and become proficient in its use. The
% system must also be efficient to ensure that employees can complete their tasks
% in a timely manner without unnecessary delays. Finally, the system must be
% error-free to ensure that employees can perform their tasks without costly mistakes
% or disruptions.

% \subsection{Secondary Stakeholders Needs}

% \subsubsection{Customers Needs}

% The main needs of the customers are to have a system that allows them to
% easily place orders or book appointments, view menus or services, and receive
% accurate information about the business. Moreover customers need easy to pay options. This need is driven by the fact that 
% customers expect a seamless and convenient experience when interacting with
% businesses, and a well-designed system can help meet these expectations.


% The main needs of the customers are to have a system that is \textbf{effective}, \textbf{efficient},
% \textbf{reliable} and \textbf{satisfying}. Customers will need to use the system
% occasionally to place orders or book appointments. Therefore the system must be
% effective to ensure that customers can easily find and use the features they
% need. The system must also be efficient to ensure that customers can complete
% their tasks quickly and without unnecessary delays. The system must be
% reliable to ensure that customers can trust the system to perform as expected
% and that their interactions with the business are not disrupted. Finally, the
% system must be satisfying to ensure that customers have a positive experience
% and are more likely to return to the business in the future.


% ------------------------
% PACT stuff below
% ------------------------



% reikia sistemizuoti kokioje situacijoje ko reikia, labiau detaliau issiaiskinti stakeholder needs.
% reikia atlikti naudojimo konteksto analize, kaip tos funkcijos isilies i vartotojo esama gyvenima. Mes galvojame apie naudojima ir konteksta vis dar.
% Naudojimo kontekstas - kaip veiklos itakoja technologiju kaita ir technologiju kaita veiklas, po to kalbesime apie pagrindines zmoniu charakteristikas, kas, kokiose veiklose ir problemas kurias norime spresti, kokiame zingsnyje kas ten stringa, kokioje aplinkoje ten vyksta, kokios tech naudojamos ir ka galime pasiulyti.
% Naudojimo kontekstas - kombinacija vartotoju, uzduociu, resursu, tikslu. Fizines, socialines, technines, kulturines, organizicines aplinkos. Jas riekes issaiskinti ir aprasyti kaip dabar veikia tas vartotojas, ka tie pasakojimai turi atspindeti: turi buti aiskus visi aspektai, kas, kokie tisklai, aplinkos, kokiais prietaisais, programomis naudojasi ir t.t.
% PS inzinerijos standartas - naudojimo kontekstas pagrindinis informacijos reikalavimu saltinis.
% Vystimasis yra begalinis ciklas, yra zmones betkokio amziaus, kutluros, jie veikia kazkokiam kontekste (requirements), tame kontekste kyla nauji reikalavimai naujoms technologijosm, naujos technologijos sudaro galimybe keistis veiklos, vel atsiras nauji reikalavimai ir t.t. gaunasi begalinis ratas.
% naujos interaktyvios tech - prisidedame dabartines zinias, sukaupta patirti ir kuriame naujus reikalavimus. Suvokdami kazkokius nepatogumus, koks yra stovis, kokias galimybes mes galime isnaudoti su technologijomis mes pasiulome sprendimus, naujas veiklas.

% People, context, activities, technologies

% ____________________________

% People - kokiomis salygomis veikia, ju fizines galimybes, darbo vieta, veikimo aplinka kur bus naudojama technologija, zmoniu ugis, aukstis, antropronetiniai dalykai? ivestis isvestis, klaviaturos ir t.t. Nera vidutinio vartotojo, universali technologija reiskia kad visiem patogu naudotis, zinoma yra tam tikri kompromisai, reikalavimai turi tilpti i budget. Pvz.: ekrano padetis, ryskumas, kontrastas, color blindness, motion sensitivity, hearing etc. etc. Musu varottoju tarpe jeigu yra vyresni zmones tai turime i tai atsizvelgti, mygtukai didesni ir t.t. Kokia turi buti darbo aplinka kad zmogus nepavargtu, nebutu broko, darbo vieta turi buti tokia kuri leistu islaikyti darbinguma.

% psihologinis erdviu suvokimas, kai kurie pasiklysta, kai kurie gerai orientuojasi erdvese. kalbos skirtumai, kai kam gali buti izeidzianti kalba, kai kam ne...

% isiminimas, trumpalaikis isiminimas, sprendimu priemimas, komunikacija, paieska - mes turime sias veiklas kurioms kuriame technologijas.

% socialiniai skirtumai - itakoja motyva pirkti nauja technologija, gali atsirasti stiprus motyvas, pvz bendrauti nuotoliniu budu, todel zmones ismoks tai. skirstosi tipai, begginer, inermediate, expert zmoniu. pvz bileteliu popieriniu nenupirksi, tai beggineriai turi tai naudoti... vidutiniskai patyre kai reikia tia naudoja o taip tai nenaudoja. profesionalai - kiekviena diena naudoja. beggineriu ir expertu poreikia yra labai skirtingi.

% kad pradedantysis galetu naudotis technologija tai ji turi vesti uz rankos, pvz vedlys, jis vedamas uz rankos.

% ekspertai - jiem reikia greicio, juos vedlys erzins.

% vidutiniskai patyre - reikia prisiminti greitai.

% jie visi turi skirtingus poreikius, reikia suprasti su kuo turime reikala.


% Mental model - mintinis modelis, koki vaizda suformavo vartotojas savo galvoje, mes bandome tai sufprasti, tai suvokimas kas ir kaip gali naudoti technologijas. Mintinis modelis nepilnas, nestabilus, mes negalime visko prisiminti, kazkas isilaiko ilgalaikeje atmintyje, kazkas dingsta. Analizuojame o ka dabar naudoja vartotojai, analizuojame ju mintines zinias. Mintinis modelis kuriamas saveikaujant su sistemomis, su kokiomis sistemomis saveikauja toki ir turesime mintini modeli.

% Nustatymas vartotojo igudzius yra musu tikslas. kokio amziaus, lytis, fizines, issilavinimas, kulturinis, motivacines galimybes, tikslai, asmenybe.
% Projektavimo tikslai turi susieti su tais skillais. musu tikslas kad vartotojai pilnai isnaudotu ka suprojektavome, norime kad is begginer greitai pereitu i intermediate.

% Pavyzdiziui kai pirma karta paleidziame programa, gali issokti gidas, padeti beginneriams, o kam nereikia tai gali uzdaryti, uzdaro vidutiniskai patyre. Ekspertam reikia kuo greiciau dirbti, kas stabdo (rankos kelimas nuo peles iki klaviaturos, galime sakyti kad viska darytu ant klaviaturos).

% kiek tu sluoksniu design reikia? tai papildomos islaidos ir t.t.


% begineriam programa turi pasakyti ka ji gali daryti, pagrindiniai dalykai ka ji daro, begineriu nera daug bet jie yra ir pirmi naudotojai todel jiems reikia viska paaiskinti. po keliu kartu dauguma begineriu tampa intermediate.
% intermediate svarbu matyti, isnaudoti visas funkcijas kurias jie zino kad yra, gebeti viska surasti, pazengusias funkcijas. GUI visas pagrindines ir pazengusias funkcijas rodo ir galima greitai surasti
% Ekspertu nera daug, kazkam galbut greiciau reikia.

% kuri viekla daznesne, ta turi buti arciau ekrano.


% Universal usability - prieinamas visiems, ir neigaliems. Panaudojamumas + prieinamumas = universalumas. kas aktualu, kam kuriame sistema, i sita turime atsakyti.



% Reikia pamineti demographics, age, occupation, gender (jei reikia), disabilities. Motyvacija ismokti technologijas, jomis naudotis. Naudojamos technologijos ir prietaisai, IT lygmuo, skillai.

% ____________________________

% Veiklos - reikia ne specifines mineti, o laiko apsketu daznas ar retas, bendradariavimo aspektus, individualus ar organizaciniai, tos veiklos sudetingumas, pasekmes, ar kritines, kokio tipo turinys toje veikloje.

% Daznis jos pirma charakteristika, jos trukme antra, laiko spaudimas, ar veikia
% ramybes busenoje ar skubos. Ar tai viena atomine veikla ar zingsniai yra, ar
% testine is zingsniu tai kazkuriame zingsnyje vartotojas gali sustoti. atsako
% laikas. Bendradarbiavimas, vienas ar grupeje, kaip zino kas ka padare, kaip
% koordinuoja, komunikuoja. Sudetingumas veiklos, ar veikla apibrezta ar ne,
% reikia suteikti vartotojui galimybe narsyti ivairias veiklas, kur vartotojas
% nori narsyti, kaip vartotojas supras ar nutrauke zingsni ir t.t. tai
% sudetingumas. kaip supras kad klaida padare, ar tai aktualu. kokie duomenys
% vaiksto, kiek ivesti ten reikia, kokia ten isvestis dabar ir kas siuloma, turi
% buti realiu laiku atnaujinamas turinys, kad nebutu blogai pavaizduotas, blogai
% atnaujintas.

% Ne pacios veiklos idomios bet ju charakteristikos.

% ____________________________

% Aplinka - fizine, kur vyksta ta veikla, patalpoje, uz patalpos, jeigu fizine aplinka lauke, tai skirtinga temperatura, kaip vartotojas su pirstinemis ar per lietu gali naudotis, apsvietimas, triuksmas ir t.t.

% Socialine aplinka - individuali, o gal grupine, ar privatumas aktualus, ar visi vienodu teisiu, ar yra super adminas.

%  Organizacinis kontekstas - teises ir t.t.

% ____________________________

% Technologijos - ka dabar naudoja zmones, kokias technologijas, ka jie dabar naudoja, tia itakoja ju mintini modeli, kaip jie dabar iveda informacija?

% Kokias technologijas naudosime - kompiuteris, programele telefone, ar t.t.

% Turime susivokti kokiu techniniu ir sistemu reikia. Kas yra geras turinys ir kokios charakteristikos galetu tobulinti tai, balsu ivedimas, AI ar whatever.

% kokia informacija ir funkcijos reikalingos sistemai, kas tures buti zinoma norinciam naudotis sitema.


% Interaktyvus produktas turi atitikti tai ko nori zmones, jis negali sugriauti kontekto, jis turi tai pagerinti.

% musu uzdavinys visus siuos dalykus issiaiskinti, parasyti kiekvienai pirminei ir antrinei vartotoju grupei aprasyti siuos dalykus ir po to dar kitai savaitei galime pradeti rasyti poreikius kas isto isplaukia ir kokiu funckiju riekia

% \sectionnonum{Results and conclusions}
% For details of what needs to be written in this section, please refer to the methodology requirements of the respective programme.


\section{PACT analysis}

\begin{center}
    \textcolor{red}{\textbf{(rough draft, dont look too deep into it yet)}}
\end{center}


We will use the \textbf{PACT} framework to analyze the needs of our system
users. That means we will look into the people and their physical and mental
capabilities, the characteristics of the activities they perform with current
systems, context in which they perform these activities, and the technologies
they use and which can improve their experience.


\subsection{Primary Stakeholders PACT analysis}

\subsubsection{Client Business Owners PACT analysis}

Lets start with analyzing the \textbf{people} aspect of PACT. 

\begin{itemize}
    \item \textbf{Physical capabilities:} Client business owners can have
    different physical capabilities, some of them might have visual or hearing
    impairments. This means that our system must have different contrast modes,
    font sizes, and support for screen readers.
    \item \textbf{Mental capabilities:} Client business owners can have
    different levels of technical expertise, mostly ranging from intermediate to
    expert. We expect that they will be able to use sales tracking, inventory
    management, and report generation features without much difficulty.
    \item \textbf{Psychological capabilities:} Client business owners can be
    from various cultural backgrounds and have different language preferences.
    This means that our system must support multiple languages and be culturally
    neutral.
\end{itemize}

Now lets analyze the \textbf{activities} aspect of PACT.

\begin{itemize}
    \item \textbf{Temporal aspects:} Client business owners will use the system
    daily, this means that the tasks like generating reports, tracking sales
    will be easy to memorize. Furthermore these tasks should be efficient to
    perform to avoid wasting time. Tasks like inventory management must be
    error-free as they can place time pressure on the user during busy periods.
    Furthermore this activity must be seperated into clear steps to avoid
    confusion. It is also necessary to have real time update of inventory data
    to avoid mistakes.
    \item \textbf{Cooperation:} Client business owners will mostly use the
    system individually, but they might need to share information with employees.
    There should be a clear way to share information.
    
    \item[] 
    \begin{center}
        \textcolor{red}{\textbf{(think deeper here, there is quite a lot of details to add)}}
    \end{center}


\end{itemize}

Now lets analyze the \textbf{context} aspect of PACT.

\begin{itemize}
    \item \textbf{Physical context:} Client business owners will use the system
    in a variety of physical contexts, including in their office, at home, or
    on the go. This means that our system should work well on different screen
    sizes and be responsive to different input methods (mouse, touch, keyboard).
    \item \textbf{Social context:} Client business owners will need to have administrative
    rights to manage employees and business operations. There should be a clear
    way to manage user roles and permissions.
\end{itemize}

Now lets analyze the \textbf{technologies} aspect of PACT.

\begin{itemize}
    \item \textbf{Technical capabilities:} Client business owners will have
    access to various devices (desktops, laptops, tablets, smartphones) and
    operating systems (Windows, macOS, iOS, Android). Our system must be
    compatible with these platforms and provide a consistent user experience.
    \item \textbf{Software capabilities:} Client business owners may use
    different software applications (CRM, ERP, Office Suite) alongside our
    system. We should consider integration options and data exchange formats
    (APIs, CSV, etc.) to facilitate smooth workflows.
    \item \textbf{Network capabilities:} Client business owners will rely on
    internet connectivity to access our system. We must ensure that our system
    performs well under different network conditions (Wi-Fi, mobile data).
\end{itemize}

\subsubsection{Employees PACT analysis}

Lets start with analyzing the \textbf{people} aspect of PACT. 

Now lets analyze the \textbf{activities} aspect of PACT.

Now lets analyze the \textbf{context} aspect of PACT.

Now lets analyze the \textbf{technologies} aspect of PACT.



\subsection{Secondary Stakeholders PACT analysis}

\subsubsection{Customers PACT analysis}

Lets start with analyzing the \textbf{people} aspect of PACT. 

Now lets analyze the \textbf{activities} aspect of PACT.

Now lets analyze the \textbf{context} aspect of PACT.

Now lets analyze the \textbf{technologies} aspect of PACT.

\printbibliography[title = {References and sources}]

% \appendix
% \renewcommand{\thesection}{Appendix \arabic{section}. }

% \section{\phantom{Appendix} Examples of citations}
% In the document \textit{bibliography.bib}, you need to add all the cited sources and after using the function \textit{\{\textbackslash cite\{name of the cited object\}\}} the corresponding source will be added to the list of literature sources.


% \textit{bibliography.bib} provides examples of some of the most commonly cited types of sources:
% \begin{itemize}
%     \item web pages (\textit{@online}) \cite{PvzInternetinisPuslapis},
%     \item datasets (\textit{@dataset}) \cite{dataset}
%     \item articles (\textit{@article}) \cite{PvzStraipsnLt, PvzStraipsnEn}, 
%     \item articles from conferences (\textit{@inproceedings}) \cite{PvzKonfLt, PvzKonfEn}, 
%     \item books (\textit{@book}) \cite{PvzKnygLt, PvzKnygEn}, 
%     \item theses (\textit{@thesis or mastersthesis/phdthesis} \cite{PvzMagistrLt, PvzPhdEn})
%     \item electronic publications (\textit{@misc}) \cite{PvzElPubLt, PvzElPubEn}
% \end{itemize}

% Examples are also provided for ChatGPT citation, both in general \cite{chatgpt_bendrai} and for a specific conversation \cite{chatgpt_pokalbis}.

\end{document}
